\chapter{Consequences of the Exact Symbol: Explicit Nash Constant and Poincaré Inequality}

This chapter mirrors \texttt{FracLap/Tools/NashConsequences.lean}.

\begin{theorem}\label{thm:nash_explicit}
\lean{FracLap.gagliardoSymbol_one_eq, FracLap.fractional_nash_explicit}
\uses{thm:exact_symbol, thm:fractional_nash}
\leanok
For $0 < \alpha < 2$, $\Lambda_\alpha(1) = (2\pi)^\alpha$, and for $f \in L^1 \cap L^2$ with finite energy,
\[
  \bigl(\lVert f \rVert_2^2\bigr)^{1+\alpha} \le \frac{3^{1+\alpha}}{(2\pi)^\alpha}\, E_\alpha(f)\, \bigl(\lVert f \rVert_1^2\bigr)^\alpha.
\]
\end{theorem}

\begin{proof}
\uses{thm:exact_symbol, thm:fractional_nash}
\leanok
Insert $\Lambda_\alpha(1) = (2\pi)^\alpha$ from Theorem~\ref{thm:exact_symbol} into Theorem~\ref{thm:fractional_nash}.
\end{proof}

\begin{lemma}\label{lem:cs_lintegral}
\lean{FracLap.lintegral_rpow_two_le}
\uses{thm:std_holder}
\leanok
For nonnegative measurable $g$, $(\int g\, \mathrm{d}\mu)^2 \le \mu(X) \int g^2\, \mathrm{d}\mu$.
\end{lemma}

\begin{proof}
\uses{thm:std_holder}
\leanok
Apply Theorem~\ref{thm:std_holder} with $p = q = 2$ to $g$ and $1$.
\end{proof}

\begin{theorem}[Fractional Friedrichs--Poincaré inequality]\label{thm:poincare_friedrichs}
\lean{FracLap.poincare_friedrichs}
\uses{lem:nash_cut, thm:exact_symbol, lem:cs_lintegral}
\leanok
Let $0 < \alpha < 2$ and let $f \in L^2(\mathbb{R}; \mathbb{C})$ vanish almost everywhere outside a measurable set $\Omega$ of finite measure. Then $\lVert f \rVert_2^2 \le 2 (2 \lvert \Omega \rvert / \pi)^\alpha E_\alpha(f)$.
\end{theorem}

\begin{proof}
\uses{lem:nash_cut, thm:exact_symbol, lem:cs_lintegral}
\leanok
By Lemma~\ref{lem:cs_lintegral}, $\lVert f \rVert_1^2 \le \lvert \Omega \rvert \lVert f \rVert_2^2$, so $f \in L^1$. Take $R = 1/(4\lvert \Omega \rvert)$ in Lemma~\ref{lem:nash_cut}: the low-frequency term is at most $2R\lvert \Omega \rvert \lVert f \rVert_2^2 = \frac12 \lVert f \rVert_2^2$ and can be absorbed, leaving $\lVert f \rVert_2^2 \le 2 (R^\alpha (2\pi)^\alpha)^{-1} E_\alpha(f) = 2(2\lvert \Omega \rvert/\pi)^\alpha E_\alpha(f)$.
\end{proof}
